\section{Local geometry and torsion detection}\label{tor:local}

The arguments in this section apply to every finite crystallographic root system. We work over $\C$ and allow the field $F\supset\C$ to vary. This extra quantifier is necessary when a slice is replaced by a closed subslice. In any one argument, the strict henselizations of all subslices are formed using the same separable closure $F^{\mathrm{sep}}$ of $F$.

A singular label $u$ of $X_w$ means $X_u\subseteq\Sing(X_w)$; it is nonmaximal when $X_u$ is not an irreducible component of this locus. This terminology always concerns singularity of the ambient variety $X_w$ along $X_u$.

\subsection{Slices and shared local rings}\label{tor:charts}
Write $C_z=BzB/B$, $C^u=B^-uB/B$, and $N_{u,z}=X_z\cap C^u$ for $u\le z$. Intersections have their actual scheme structures. The slice has distinguished vertex $uB/B$ and dimension $\ell(z)-\ell(u)$. We use the multiplication chart
\[
 O_u\cap X_z\simeq C_u\times N_{u,z},
 \qquad O_u=u(B^-B/B).
\]
The same unipotent multiplication constructs these charts for every $z$, so they respect inclusions; see \cite[Lemma~2.1 and its proof]{KWY13}. Schubert varieties are normal and Cohen--Macaulay, and therefore their slices have these properties as well \cite[Sections~2.1--2.2]{Bri05}.

For $F\supset\C$, set
\[
 \begin{gathered}
 R_F(u,z)=F[N_{u,z}]_{\mathfrak m_u},
 \qquad A_F(u,z)=R_F(u,z)^{\sh},\\
 U_F(u,z)=\Spec A_F(u,z)\setminus\{s\},
 \end{gathered}
\]
where $\mathfrak m_u$ is the vertex ideal and $s$ is the closed point. The rings $R_F(u,z)$ and $A_F(u,z)$ are Noetherian excellent normal local domains, and they are Cohen--Macaulay. Indeed, over the algebraically closed field $\C$, normal integral varieties are geometrically normal and geometrically integral. Normality persists under extension of the ground field \cite[Tags~038N and 038P]{The26}; the Cohen--Macaulay property does so here because the algebras are of finite type \cite[Tag~045P]{The26}. Strict henselization preserves the required local properties \cite[Tags~06DI and 06LM]{The26}.

Fix $u\le z\le w$, and let $I_z$ be the ideal of the closed subslice. Compatibility of strict henselization with quotients identifies
\[
 A_F(u,z)=A_F(u,w)/I_zA_F(u,w)
\]
with the chosen common separable closure \cite[Tag~05WS]{The26}. Consequently the $U_F(u,z)$ are closed subschemes of a common punctured spectrum. Each is nonempty and connected when $z>u$.

The product chart also gives
\[
 \OO_{X_w,\eta_{X_u}}
 \simeq\bigl(\C(C_u)[N_{u,w}]\bigr)_{\mathfrak m_u}.
\]
Thus the ordinary generic-point ring is obtained by taking $F=\C(C_u)$; the smaller subslices in the same argument must still use this field.

\subsection{Reduction to an effective Levi subgroup}
\begin{proposition}\label{tor:levi}
If $N=N_{u,w}$ has dimension $d$, it is isomorphic, as an affine scheme with vertex and ambient torus action, to a Schubert slice for a Levi subgroup whose derived rank is at most $d$. The isomorphism preserves root lengths after conjugating that Levi to a standard Levi subgroup, and it commutes with every extension $F\supset\C$.
\end{proposition}
\begin{proof}
Restrict the root-coordinate functions of the opposite cell to $N$, discarding those that vanish identically. They generate $\C[N]$ and have root characters $\alpha_i$. Let $V$ be their rational span, of dimension $r$. Integrality gives a point at which all these functions are nonzero. Its stabilizer is their common kernel, so its orbit has dimension $r$; in particular $r\le d$.

Let $H$ be the identity component of that kernel. A character is trivial on $H$ precisely when it belongs to $V$ over $\Q$. Thus the roots of $G'=Z_G(H)$ are $\Phi\cap V$, and they span $V$ because they include the $\alpha_i$. Choose a cocharacter $\lambda$ of $H$ avoiding the finitely many root hyperplanes that do not contain $H$. Then $G'=Z_G(\lambda)$ is a Levi subgroup of derived rank $r$. If $H$ is trivial, take $\lambda=0$ and $G'=G$; the same character calculation then gives $r=\rank G$. Since its coordinate generators have weight zero under $\lambda$, $N$ is fixed as a scheme.

Let $S$ be the connected component of $(G/B)^\lambda$ containing $N$. The pattern map gives a $G'$-equivariant isomorphism $\psi:S\longrightarrow G'/B'$, where $B'=B\cap G'$. For every fixed-point label $v$ in this component, it identifies the fixed ordinary cell $(C_v)^\lambda$ and the fixed opposite cell $(C^v)^\lambda$ with $C'_{\phi(v)}$ and $(C')^{\phi(v)}$, respectively \cite[Theorem~10]{BB03}\cite[Theorem~4.1 and the following paragraph]{Woo10}. These are statements about the fixed cells themselves; no irreducibility assertion for $X_w\cap S$ is needed.

The open Richardson stratum $C_w\cap C^u$ is nonempty and dense in $N$. Since it is fixed, the connected fixed cell $(C_w)^\lambda$ meets $S$, and hence lies in $S$; in particular $w\in S$. Likewise $(C^u)^\lambda\subset S$. The two cell isomorphisms are restrictions of the same injective map $\psi$, so they identify their intersection with the entire intersection of their images:
\[
 \psi(C_w\cap C^u)=C'_{\phi(w)}\cap(C')^{\phi(u)}.
\]
Take reduced closures inside the fixed opposite cell on the left and inside $(C')^{\phi(u)}$ on the right. The first closure is $N$, since $N$ is reduced and closed in $C^u$, and the indicated stratum is dense. The second is the Levi Schubert slice $N'_{\phi(u),\phi(w)}$. This gives an isomorphism of actual affine schemes. The same argument applies to every $N_{u,z}\subset N$ with $u\le z\le w$, so the identifications respect closed subslices.

Since $T\subset G'$ and $\psi$ is $G'$-equivariant, it preserves the ambient characters. Conjugation to a standard Levi preserves root lengths. All maps were constructed over $\C$, so the isomorphism persists after any extension of the ground field.
\end{proof}

\subsection{Finite covers and the Schubert boundary}
\begin{lemma}\label{tor:closed-gluing}
Suppose finitely many closed subschemes $T_a\subset T$ cover its support, and $P\longrightarrow T$ is finite \'etale. Sections on the $T_a$ that agree on their schematic pairwise intersections glue uniquely to a section on $T$. It suffices to check agreement on a family covering the support of each intersection.
\end{lemma}
\begin{proof}
The chosen sheet and the union of the other sheets above each $T_a$ are closed in $P$. Agreement makes the finite unions of these two families disjoint, and together they cover $P$. The chosen union is therefore clopen. With its induced scheme structure it is finite \'etale over $T$ and has one point in each geometric fibre, hence is isomorphic to $T$. Finally, the equalizer of two sections of a finite \'etale map is clopen. Agreement on its whole support is therefore sufficient, even if the intersection is nonreduced.
\end{proof}

\begin{theorem}\label{tor:boundary}
Let $d=\ell(w)-\ell(u)\ge4$. A finite \'etale cover of $U_F(u,w)$ that splits on every $U_F(u,v)$ with $u<v\lessdot w$ is split. Moreover, restriction induces an injection
\begin{equation}\label{tor:boundary-injection}
 \Tor\Pic U_F(u,w)\longrightarrow
 \bigoplus_{u<v\lessdot w}\Tor\Pic U_F(u,v).
\end{equation}
\end{theorem}
\begin{proof}
On the simply connected form of the group, a boundary section of $\mathcal L_\rho$ defines an effective Cartier divisor supported on all Schubert divisors, with positive multiplicities \cite{AS14}. The flag variety is unchanged by this choice of isogeny. Restrict the section to $N_{u,w}$; it is nonzero because the slice meets the top cell. Trivializing at the vertex gives a nonzero element $f$ of the maximal ideal of $A=A_F(u,w)$. Write $U=\Spec A\setminus\{s\}$ and $T_0=\Spec(A/fA)\setminus\{s\}$. Its support is covered by the closed coatom pieces $U_v$.

The support of an intersection of lower pieces $U_{z_i}$ is covered by the common lower pieces $U_q$ with $u<q\le z_i$ for every $i$. This follows by restricting the Schubert support intersection identity to the opposite cell, using the common field and the quotient compatibility above. The case $q=u$ disappears on the puncture. No reducedness assertion about these intersections is needed.

The order complex of the open Bruhat interval $(u,w)$ is a sphere $S^{d-2}$ \cite{BW82}; in this range it is connected and simply connected. Because every $U_z$ with $u<z<w$ is connected, a cover of $T_0$ split on each coatom is split on every lower piece. The section sets are finite, and the restriction maps are bijections satisfying the usual composition identities. Transport labels along comparability edges. Backtracking gives the identity, and each triangle gives the identity by composition. Simple connectivity therefore makes transport independent of the path, and all these finite sets can be labelled compatibly.

Sections with the same label agree on every common lower piece, hence on the entire support of each schematic intersection. Lemma~\ref{tor:closed-gluing} glues them on $T_0$, so the restricted cover splits there. The ring $A$ is henselian and has depth $d$. By \cite[Tag~0BLJ]{The26}, restriction from finite \'etale covers of $U$ to those of $T_0$ is fully faithful. Lift a splitting isomorphism and its inverse; faithfulness makes the lifts inverse. This proves the first assertion. Essential surjectivity is not used.

For the second assertion, take $L$ in the kernel and choose $L^{\otimes m}\simeq\OO_U$. Its Kummer torsor restricts to a split torsor on each $U_v$: since $L|_{U_v}$ is trivial, the restricted torsor class can only come from a unit modulo $m$th powers. Normal Hartogs extension identifies these units with $A_F(u,v)^*$, and every such unit has an $m$th root by strict henselianity in characteristic zero. The first assertion splits the original torsor, so $L$ is trivial.
\end{proof}

\subsection{An integral curve-degree bound}
\begin{theorem}\label{tor:degrees}
Let $A$ be a Noetherian strictly henselian normal local $\Q$-algebra of dimension at least two, and put $U=\Spec A\setminus\{s\}$. Let $f:Y\longrightarrow\Spec A$ be proper birational, with $Y$ regular and integral. Suppose every finite \'etale cover of its connected closed fibre $F_0$ splits. Let $E_1,\ldots,E_r$ be all prime divisors of $Y$ contained in $F_0$. For actual proper integral $1$-cycles $Z_i$ on $F_0$, allowing signed coefficients, set
\[
 N_{ij}=\deg(\OO_Y(E_j)|_{Z_i}).
\]
If $N$ has full column rank and largest Smith invariant $h$, then $h\Tor\Pic(U)=0$. For $r=0$, take $h=1$.
\end{theorem}
\begin{proof}
Choose $L^{\otimes m}\simeq\OO_U$. Its generic Kummer algebra is $K[z]/(z^m-a)$ for $a\in K^*$, where $K=\Frac A$; this algebra need not be a field. The pulled-back torsor is \'etale on $Y\setminus F_0$. At every prime divisor not contained in $F_0$, the DVR description of a Kummer torsor shows that the valuation of $a$ is divisible by $m$. Consequently
\[
 \operatorname{div}_Y(a)=mD+\sum_jc_jE_j
\]
for an integral divisor $D$ and integers $c_j$. All these divisors are Cartier because $Y$ is regular.

Taking degrees on the $Z_i$ gives $Nc\equiv0\pmod m$. Smith normal form supplies an integer matrix $M$ with $MN=hI_r$. Hence $hc=mb$ for an integer vector $b$, and
\[
 \operatorname{div}_Y(a^h)=mD',
 \qquad D'=hD+\sum_jb_jE_j.
\]
Locally write $D'=\operatorname{div}(g_\alpha)$. The algebras
\[
 \OO_Y[t_\alpha]/(t_\alpha^m-a^h/g_\alpha^m),
 \qquad t_\alpha=(g_\beta/g_\alpha)t_\beta,
\]
glue by the indicated transitions to a finite \'etale Kummer torsor on $Y$. Proper henselian equivalence \cite[Tag~0A48]{The26} and the hypothesis on $F_0$ split this torsor. Thus $a^h$ is an $m$th power in $K$. The $h$th power of the original torsor is generically split, hence split on the normal scheme $U$. In particular $L^{\otimes h}$ is trivial.
\end{proof}

Only genuine curve degrees and the complete set of columns are required. The rows need not generate all curve degrees. In particular, a full-column minor equal to $\pm1$, or an integer combination of full-column maximal minors equal to $1$, proves vanishing of torsion.

\subsection{Bott--Samelson sources and exact exceptional columns}\label{tor:sources}
Fix a reduced word $Q=(i_1,\ldots,i_l)$ for $w$, with Bott--Samelson resolution $f_Q:Y_Q\longrightarrow X_w$ and canonical boundaries $D_1,\ldots,D_l$. The cell's unipotent section and $B$-equivariance give the compatible product
\[
 f_Q^{-1}(O_u\cap X_w)\simeq C_u\times Y_N,
 \qquad Y_N=f_Q^{-1}(N_{u,w}).
\]
The product is obtained by acting on $Y_N$ with the unique cell coordinate in the chosen unipotent section, and it transports the boundary bundles. Its left side is smooth and integral, so $Y_N$ is also smooth and integral; the map $Y_N\longrightarrow N_{u,w}$ is proper birational.

After extension to $F$ and base change to $A=A_F(u,w)$, the resulting $Y_A$ is regular, since localization and ind-\'etale base change preserve regularity. It is integral as well. Indeed, proper birationality to a normal target gives $f_*\OO=\OO$; flat base change preserves this identity \cite[Tag~02KH]{The26}, so $Y_A$ is connected. Distinct irreducible components of a regular Noetherian scheme are disjoint and open, which proves integrality. Its generic fibre remains a point.

\Needspace{5\baselineskip}
\begin{lemma}\label{tor:fibre}
The closed fibre of $Y_A$ is connected, and every finite \'etale cover of it splits.
\end{lemma}
\begin{proof}
The complex Bott--Samelson fibre $F_u=f_Q^{-1}(u)$ is connected by normality of the target. It is affine-paved: its intersections with the projective Bia\l ynicki--Birula cells are the polynomial graphs of \cite[Proposition~2.1]{Ha04}, and the closed Bia\l ynicki--Birula filtration restricts to the fibre.

For a connected proper complex affine-paved scheme, compactly supported cohomology of the paving has only even degrees. The number of zero-dimensional cells is therefore its zeroth Betti number, namely one. A finite \'etale cover of degree $n$ pulls each affine cell back to $n$ copies, because complex affine space has no nontrivial finite \'etale cover by Riemann existence \cite[Expos\'e~XII, Theorem~5.1]{Gro71}. The cover is proper and has $n$ zero-dimensional cells. It consequently has $n$ connected components, each of degree one, and is split.

Nilpotents do not affect the finite \'etale category. The closed fibre over $A$ is $F_u\times_{\C}F^{\mathrm{sep}}$. In characteristic zero, $F^{\mathrm{sep}}$ is algebraically closed. Proper invariance of the finite \'etale category under extension of algebraically closed fields \cite[Tag~0A49]{The26} proves the assertion.
\end{proof}

\begin{lemma}\label{tor:columns}
Set $z_j=\Dem(Q\setminus j)$. The prime divisors of $Y_A$ contained in its closed fibre are precisely the restrictions of the $D_j$ with $z_j=u$, each reduced and with multiplicity one. In particular the exact column set is $E(Q,u)=\{j:\Dem(Q\setminus j)=u\}$.
\end{lemma}
\begin{proof}
The divisor $D_j$ is the smooth integral Bott--Samelson variety for the deleted word, whether or not that word is reduced; its image is $X_{z_j}$. If $u\le z_j$, its restriction to the source product of Subsection~\ref{tor:sources} has the form $C_u\times D_{j,N}$. If $z_j=u$, the smooth integral $D_{j,N}$ lies entirely over the vertex. It is geometrically integral, so extension to $F$ followed by the common strict base change leaves a single reduced prime divisor. Its bundle is the restriction of $\OO_{Y_Q}(D_j)$.

For completeness, first work over $F$, rather than comparing heights across a possibly transcendental extension of $\C$. A height-one closed-fibre prime on $Y_A$ contracts, by flat going down, to a prime of height at most one on $Y_{N,F}$. It cannot contract to the generic point, whose image is generic in $N_F$ rather than the vertex. It thus contracts to a prime divisor over the vertex. Its product with $C_{u,F}$ closes to a prime divisor on $Y_{Q,F}$ outside the inverse image of the top cell. That complement is the union of the canonical $D_{j,F}$, so this prime is one of them. Its image centre on the slice is $N_{u,z_j}$, which lies over the vertex exactly when $z_j=u$. Distinct columns have distinct contractions, so they cannot merge after strict base change. This proves completeness and the stated multiplicities.
\end{proof}

\begin{lemma}\label{tor:products}
Let $N=N_1\times\cdots\times N_s$, where $s\ge2$ and every $N_i$ is a positive-dimensional actual Schubert slice. Then, for every $F\supset\C$, the strict puncture at the product vertex has torsion-free Picard group. The same conclusion holds if the factors are scalar Veronese cones.
\end{lemma}
\begin{proof}
Take the product of the regular proper birational slice sources just constructed. The closed fibre of each factor is a proper closed subset of an integral source, hence has codimension at least one. Their product has codimension at least two and contains no prime divisor. It is connected and affine-paved, so its finite \'etale covers split by the preceding argument.

After extension to $F$ and strict base change, regularity and integrality persist as above. A new height-one closed-fibre prime would contract by flat going down to a prime of height at most one in the original product source, contradicting the codimension of its vertex fibre. The empty-column case of Theorem~\ref{tor:degrees} applies. For a scalar Veronese cone of dimension $m\ge2$, use the regular source $\operatorname{Tot}(\OO_{\PP^{m-1}}(-2))$; a one-dimensional factor is smooth and uses the identity. The same product argument has projective spaces or points as its closed fibres.
\end{proof}

\subsection{Mask characters as integral curve degrees}
Use the adjoint form for this calculation, so that $X^*(T)$ is the root lattice and every root is primitive. This change leaves the flag variety and the Bott--Samelson variety unchanged; it is independent of the simply connected form used for $\mathcal L_\rho$ above.

For a mask $\gamma\in\{0,1\}^l$, write $\gamma^j=s_{i_1}^{\gamma_1}\cdots s_{i_j}^{\gamma_j}$. Its endpoint is $\gamma^l$. The canonical linearization on $\OO(D_j)$ has fibre character
\begin{equation}\label{tor:characters}
 \chi_j(\gamma)=
 \begin{cases}
  0,&\gamma_j=1,\\
  -\gamma^{j-1}\alpha_{i_j},&\gamma_j=0.
 \end{cases}
\end{equation}
Indeed, the invariant boundary section is nonzero at a selected position; at an omitted position its normal character is the transported negative simple root. The bundle is pulled back through later stages of the tower, so later coordinates do not change this character.

\begin{proposition}\label{tor:mask-degrees}
Fix an endpoint-$u$ reference mask $\gamma_0$ and simple-root coordinate functionals $\lambda_a$ on the root lattice. Every row
\begin{equation}\label{tor:raw-rows}
 \bigl(\lambda_a(\chi_j(\gamma)-\chi_j(\gamma_0))\bigr)_{j\in E(Q,u)},
 \qquad \gamma^l=u,
\end{equation}
is the simultaneous degree row of a signed integral proper $1$-cycle on the closed fibre. Consequently, a full-column-rank matrix of these rows with largest Smith invariant $h$ annihilates $\Tor\Pic U_F(u,w)$ for every $F\supset\C$.
\end{proposition}
\begin{proof}
First join endpoint-$u$ masks by moves changing two bits and keeping the endpoint. Coxeter deletion successively removes two selected letters from a nonreduced selected expression, ending at a reduced $u$-subword. These reduced subwords are complements of the facets of the subword complex $\Delta(Q,u)$. Its pure shellability \cite[Theorems~2.5 and 3.7]{KM04} makes its facet adjacency graph connected; an adjacency exchanges two positions. Thus the entire endpoint-mask graph is connected.

For one such move, let $p<q$ be its changed positions, and let $\beta$ be the positive representative of $\pm\gamma^{p-1}\alpha_{i_p}$. Direct multiplication gives
\[
 \delta^t=
 \begin{cases}
  \gamma^t,&t<p\text{ or }t\ge q,\\
  s_\beta\gamma^t,&p\le t<q.
 \end{cases}
\]
Equality of the endpoints forces $\gamma^{q-1}s_{i_q}(\gamma^{q-1})^{-1}=s_\beta$. Thus both changed positions are $\beta$-walls. The masks agree at every other position and are equivalent for the codimension-one subtorus $\ker\beta$ in the terminology of \cite[Sections~1 and 4]{Ha04}.

Here is the precise fibre geometry of this equivalence class. List all its $\beta$-walls as $a_1<\cdots<a_m$, and let $G_\beta$ be the rank-one subgroup generated by $T$ and the root subgroups for $\pm\beta$. After choosing the minimal gallery in the class, the partial-gallery parametrization of \cite[equation~(4.1), Proposition~4.2 and Corollary~4.4]{Ha04} identifies its fixed component with $(G_\beta/B_\beta)^m\simeq(\PP^1)^m$. Explicitly, between walls $a_j$ and $a_{j+1}$, the partial flag is $g_j\gamma_{min}^tB$, with $g_jB_\beta$ the $j$th coordinate. The flags at the walls recover these coordinates, giving an isomorphism of schemes, not merely a parametrization of closed points.

The final flag is $g_m\gamma_{min}^lB$. Hence the endpoint map is the last projection followed by the isomorphism onto the rank-one orbit in $G/B$. Its fibre over $u$ is exactly $(\PP^1)^{m-1}\times\{u\}$. In particular $m\ge2$ for the two distinct masks under consideration. Successively changing the first $m-1$ coordinates joins these masks by coordinate-axis projective lines in the actual fibre $f_Q^{-1}(u)$.

On each such line, the tangent character at a fixed endpoint is $\beta$ or $-\beta$. For a linearized line bundle, its two fibre characters differ by its degree times the negative of the tangent character at the starting endpoint. This follows from the transition function for $\OO_{\PP^1}(n)$ on the two standard affine charts. Thus, for an oriented axis path $C_1,\ldots,C_t$, there are signs $\eps_\nu\in\{1,-1\}$, independent of the line bundle, such that the signed integral cycle $Z_e=\sum_\nu\eps_\nu[C_\nu]$ satisfies
\[
 \chi_j(\delta)-\chi_j(\gamma)
 =\beta\deg(\OO(D_j)|_{Z_e})
\]
simultaneously for all $j$. Along a path from $\gamma_0$ to $\gamma$, apply $\lambda_a$ and add. This yields \eqref{tor:raw-rows}, realized by the cycle $\sum_e\lambda_a(\beta_e)Z_e$. Every coefficient is integral. The cycles lie over $u$, hence in the slice source, and the product chart transports their boundary bundles. Degrees are unchanged after extension to the algebraically closed residue field. Apply Theorem~\ref{tor:degrees} and Lemmas~\ref{tor:fibre}--\ref{tor:columns}.
\end{proof}

The full raw row lattice also equals the lattice of the divided two-bit-move rows $((\chi_j(\delta)-\chi_j(\gamma))/\beta)_j$. One inclusion follows from the path argument. Conversely, a primitive root $\beta$ has relatively prime simple-root coordinates, so an integral B\'ezout combination of its coordinate differences recovers the divided row. Edge differences are differences of reference rows. This equality survives projection to a fixed column set. Rational normalization of rows, or replacement by a different ambient lattice, would not preserve the argument.

\paragraph{A small degree witness.}
In type $A_2$, take $Q=(1,2,1)$, $u=s_1$, and endpoint masks $\gamma=100$ and $\delta=001$. The full exceptional column set is $E(Q,u)=\{2\}$. Formula~\eqref{tor:characters} gives $\chi_2(\gamma)=-(\alpha_1+\alpha_2)$ and $\chi_2(\delta)=-\alpha_2$, so their difference is $\alpha_1$. The two $\alpha_1$-walls are positions $1$ and $3$; their fixed component is $\PP^1\times\PP^1$, and its endpoint-$u$ fibre is a coordinate line. With the sign prescribed above, this line realizes the raw simple-root rows $(1)$ and $(0)$, the latter by the zero cycle. Thus the undivided character coordinates already give integral degree rows.

\subsection{Passage to ordinary points and partial flags}
We use flat compatibility of duals and reflexive powers from Lemma~2.2.

\begin{lemma}\label{tor:strict-injection}
For an excellent normal local domain $A$ of dimension at least two, pullback to $A^{\sh}$ is injective on the Picard group of the punctured spectrum.
\end{lemma}
\begin{proof}
The closed fibre of $A\longrightarrow A^{\sh}$ is a field, so the upper puncture is the inverse image of the lower puncture. Extend a line bundle on the latter to its unique finite rank-one reflexive module $M$ on $\Spec A$: take a rational section and close its finite Weil divisor. Normality and the codimension of the complement give uniqueness. By Lemma~2.2, its pullback is reflexive. If the pulled-back line bundle is trivial, normal Hartogs uniqueness makes this module $A^{\sh}$. Faithfully flat descent makes $M$ free of rank one, proving injectivity.
\end{proof}

\begin{proposition}\label{tor:ordinary}
For a full-flag Schubert variety $X$, ordinary torsion detection at all scheme points is equivalent to vanishing of punctured Picard torsion at $\eta_Y$ for every Schubert subvariety $Y\subseteq\Sing(X)$ that is not an irreducible component of $\Sing(X)$. It suffices to prove this vanishing after strict henselization.
\end{proposition}
\begin{proof}
First suppose these punctured groups vanish. A local class at $x$ extends to a global Weil class by taking closures of its prime divisors. The affine big cell has trivial class group, so divisor localization represents the global class by a $B$-stable boundary divisor $D$. Suppose its class at $x$ is torsion, all maximal singular generic localizations through $x$ vanish, but $D$ is not Cartier at $x$.

Its closed non-Cartier locus is $B$-stable. A component $Y$ containing $x$ is therefore a Schubert subvariety. At $\eta_Y$, the class is nonzero torsion and is Cartier on the entire puncture of $\OO_{X,\eta_Y}$: a nonfree point there would give a strictly larger nonfree subvariety containing $Y$. Regular local rings are factorial, so $Y\subseteq\Sing(X)$ and $\codim_XY\ge2$. It is not a maximal singular component, because the corresponding detector was assumed zero. Reflexive Hartogs extension shows that triviality on this puncture would imply Cartierness at $\eta_Y$. Thus we have obtained a nontrivial torsion line bundle on a puncture whose Picard torsion vanishes by hypothesis, a contradiction.

Conversely, extend a torsion line bundle on the puncture at such an $\eta_Y$ reflexively. A reflexive power of this extension is trivial by the same extension argument. Every maximal singular component containing $Y$ strictly contains it, so its generic point belongs to the puncture and has zero detector. Torsion detection at $\eta_Y$ makes the extension free, proving the required vanishing. The final assertion is Lemma~\ref{tor:strict-injection}.
\end{proof}

\begin{proposition}\label{tor:partial}
If ordinary torsion detection holds at all scheme points of every full-flag Schubert variety, it also holds for every partial-flag Schubert variety.
\end{proposition}
\begin{proof}
For $X\subset G/P$, its inverse image $Y\subset G/B$ is a Schubert variety, and $\pi:Y\longrightarrow X$ is smooth, surjective and Zariski locally trivial, with geometrically irreducible fibre $P/B$. Thus $\Sing Y=\pi^{-1}(\Sing X)$, and maximal singular components correspond by inverse image. Indeed, the inverse image of an irreducible base under this locally trivial map has irreducible fibres and is irreducible; no vertical component occurs.

At an arbitrary point $x$, choose in a local trivialization the generic point $y$ of an affine big cell of its fibre. Writing $A=\OO_{X,x}$, its ring is the faithfully flat local extension
\[
 B=A[t_1,\ldots,t_a]_{\mathfrak m_AA[t_1,\ldots,t_a]}.
\]
If $\mathfrak p\subset\mathfrak m_A$ is the prime of a maximal singular component through $x$, the corresponding upstairs generic localization is
\[
 B_{\mathfrak pB}
 =A_{\mathfrak p}[t_1,\ldots,t_a]_{\mathfrak pA_{\mathfrak p}[t_1,\ldots,t_a]}.
\]
These are actual compatible detector maps. A finite rank-one reflexive module representing a downstairs torsion class remains reflexive and torsion over $B$ by Lemma~2.2. Vanishing of its downstairs detectors gives vanishing upstairs. Detection on $Y$ makes it free over $B$, and faithfully flat descent makes it free over $A$. No isomorphism of the entire class groups is required.
\end{proof}

\section{The finite-type torsion theorem}\label{tor:finite}

Retain the objects of Section~\ref{tor:local}, and abbreviate $T_F(N)=\Tor\Pic(\Spec R_F(N)^{\sh}\setminus\{s\})$. Every assertion below quantifies over all fields $F\supset\C$. We prove the local vanishing required by Proposition~\ref{tor:ordinary}. The only computational inputs are complete bounded calculations with integer certificates.

\subsection{The finite certificates and what they prove}
Choose one reduced word $Q_w$ for each element of the finite Weyl group under consideration. Testing every reduced expression is unnecessary: each pair needs only one regular source. For $u\le w$, use the full column set $E(Q_w,u)$ from Lemma~\ref{tor:columns}. If it is empty, Proposition~\ref{tor:mask-degrees} already gives $T_F(N_{u,w})=0$. Otherwise the rows are the literal integer rows of \eqref{tor:raw-rows}; no rational row normalization is allowed.

Here is the finite algorithm for the bounded inputs. Starting at the identity, breadth-first search by right multiplication in the simple reflections constructs the group and one shortest word for every element. Closure under every generator proves completeness, and the first word found is reduced. For each word, compute every single deletion by the Demazure rule: multiply by a letter exactly when it increases length. Retain the resulting exact centre $u$ and every position producing it. Evaluate masks by ordinary group multiplication and characters by \eqref{tor:characters}. Every displayed or stored minor uses every column, with distinct valid row indices. Calculate determinants and B\'ezout identities over the integers.

\begin{proposition}\label{tor:certificates}
The following bounded cases have primitive full-column row lattice:
\begin{enumerate}
 \item in $B_3$, every exact-centre pair of dimension three;
 \item in $C_3$, every exact-centre pair of dimension three except $(Q_w,Q_u)=(121321,121)$;
 \item in $F_4$, every exact-centre pair of dimension four.
\end{enumerate}
The exceptional $C_3$ pair has one column and raw difference lattice $2\Z$. In rank two, the only possible nonzero type-$C$ surface input is $(Q_w,Q_u)=(121,1)$, again with bound two.
\end{proposition}
\begin{proof}
For $B_3$ and $C_3$, use signed permutations of three coordinates, with simple roots $\alpha_1=e_1-e_2$, $\alpha_2=e_2-e_3$, and $\alpha_3=e_3$ in $B_3$ or $\alpha_3=2e_3$ in $C_3$. Each group has $48$ elements. The algorithm gives $199$ distinct exact-centre pairs in each type: $138$ have gap one, $49$ have gap two, and $12$ have gap three. Thus there are $61$ exceptional exact-centre pairs of gap at least two. Table~\ref{tor:bc-table} lists all twelve gap-three pairs. Each matrix uses all exact columns and has determinant $\pm1$, except for the indicated $C_3$ pair.

For that exceptional pair, the endpoint masks are $001011$, $011010$, $100011$, $110001$, and $111000$. Their exceptional-column characters, in simple-root coordinates, are respectively
\[
 (0,0,-1),\quad (0,-2,-1),\quad (0,0,-1),\quad
 (-2,-2,-1),\quad (-2,-2,-1).
\]
Hence all raw differences are even, and some coordinate difference is $2$ or $-2$. The complete raw difference lattice is therefore $2\Z$.

The rank-two enumeration gives the surface assertion. In the $C_2$ convention, $\alpha_1=e_1-e_2$ is short and $\alpha_2=2e_2$ is long. For $Q=121$ and $u=s_1$, the exact column is position $2$, and the endpoint masks are $100$ and $001$. The character difference in this column is $2(e_1-e_2)$, whose simple-root coordinate column is $(2,0)^{\mathsf t}$ up to overall sign. There are exactly three nonempty dimension-two exact column sets in $C_2$:
\[
\begin{array}{c c c c}
 Q_w&Q_u&E(Q_w,u)&\text{full-column degree matrix}\\[2pt]
 121&1&\{2\}&(2)\\
 212&2&\{2\}&(1)\\
 1212&12&\{2,3\}&\begin{pmatrix}0&-1\\1&-1\end{pmatrix}
\end{array}
\]
For the last pair, reference mask $0011$ and mask $1100$ give exactly the two displayed root-coordinate rows. This is also a direct rank-two check, independent of the enumeration. Every other pair of the stated dimension has empty exact column set. The group operations in $B_3$ and $C_3$ agree, but the character matrices differ because the root lengths differ.

For $F_4$, use simple roots $e_2-e_3$, $e_3-e_4$, $e_4$, and $\tfrac12(e_1-e_2-e_3-e_4)$. The integral reflection matrices in the simple-root basis generate $1152$ elements, with maximum length $24$. Every simple reflection acts within the closed enumerated set, so this is exactly the generated Weyl group. All single deletions give $567$ dimension-four exact-centre pairs: $497$ have one column, $49$ have two, and $21$ have three. The accompanying reconstructor independently generates all these pairs and their raw mask-coordinate rows. For each pair it records full-column minors and integer B\'ezout coefficients giving a combination equal to $1$.

For a dimension-four pair, an endpoint mask selects at least $\ell(w)-4$ letters and has the same parity as $\ell(w)-4$. The complete reduced word cannot have the shorter endpoint. Thus it suffices to enumerate masks omitting two or four letters. Each stored witness is also checked directly by its endpoint and characters. The supplementary calculation retains complete group and centre coverage, all exceptional columns, legal row witnesses, and exact determinant identities. A list of numerical gcds alone would not establish these facts. Proposition~\ref{tor:mask-degrees} now gives the stated consequences for the actual punctured local objects over every $F$.
\end{proof}

\begin{table}[tbp]
\centering
\small
\renewcommand{\arraystretch}{1.3}
\begin{tabular}{@{}ccccc@{}}
\toprule
$Q_w$&$Q_u$&Exact positions&$B_3$ square matrix&$C_3$ square matrix\\
\midrule
$13213$&$13$&$3$&$(-1)$&$(-1)$\\
$121321$&$121$&$4$&$(-1)$&$(-2)$\\
$132132$&$132$&$3$&$(-1)$&$(-1)$\\
$213213$&$213$&$4$&$(1)$&$(1)$\\
$1213213$&$1213$&$4,5$&$\begin{pmatrix}-1&0\\-1&1\end{pmatrix}$&$\begin{pmatrix}-2&1\\1&-1\end{pmatrix}$\\
$1321323$&$1323$&$3$&$(-1)$&$(-1)$\\
$2132132$&$2132$&$4$&$(1)$&$(1)$\\
$2321323$&$2323$&$4$&$(-1)$&$(-1)$\\
$12132132$&$12132$&$4,5$&$\begin{pmatrix}-1&0\\-1&1\end{pmatrix}$&$\begin{pmatrix}-2&1\\1&-1\end{pmatrix}$\\
$12321323$&$12323$&$5$&$(-1)$&$(-1)$\\
$21321323$&$21323$&$4$&$(1)$&$(1)$\\
$121321323$&$121323$&$4,5,6$&$\begin{pmatrix}-1&0&1\\0&0&-1\\-1&1&0\end{pmatrix}$&$\begin{pmatrix}-2&0&1\\1&0&-1\\-2&1&0\end{pmatrix}$\\
\bottomrule
\end{tabular}
\caption{Complete dimension-three exact-centre list in $B_3$ and $C_3$, with witnesses reconstructed from raw character rows. Every displayed matrix has determinant $\pm1$ except $(-2)$. Positions are one-based and comprise the entire exact column set. A digit $1$ in a word denotes a simple reflection, not the identity.}\label{tor:bc-table}
\end{table}

\Needspace{8\baselineskip}
\subsection{Simply laced types and type \texorpdfstring{$B$}{B}}
\begin{proposition}\label{tor:simply-laced-b}
For a simply laced slice, $T_F(N)=0$ in every dimension at least two. For a type-$B$ slice it vanishes in every dimension at least three.
\end{proposition}
\begin{proof}
In dimensions at most three, Proposition~\ref{tor:levi} reduces to a Levi subgroup of rank at most three. If at least two positive-dimensional factors occur, apply Lemma~\ref{tor:products}. The irreducible simply laced factors are $A_1$, $A_2$, and $A_3$. The first two have only smooth slices. In $A_3$, the only singular three-dimensional pairs are $(u,w)=(1324,3412)$ and $(2143,4231)$. Their rank equations reduce respectively to $c=0$, $ae+bd=0$, and to $bc-ad=0$, up to independent smooth variables; at dimension three these are ordinary threefold nodes.

A regular small source of the node is $\operatorname{Tot}(\OO_{\PP^1}(-1)\oplus\OO_{\PP^1}(-1))$. Its vertex fibre is $\PP^1$ and contains no divisor, so the empty-column case of Theorem~\ref{tor:degrees} kills torsion. These low-rank equations can also be read from the standard generic-singularity descriptions \cite{Man01a,Cor03,BW03}. As an independent check, the supplementary reconstructor verifies all dimension-two and dimension-three exact-centre inputs in $A_3$; their row lattices are primitive, and all omitted pairs have empty column set.

For a single type-$B$ factor, the rank-three input of Proposition~\ref{tor:certificates} gives the dimension-three base. Smaller standard Levi factors are included: restricting the signed-permutation computation to a smaller alphabet gives the complete rank-one or rank-two computation. Empty column sets are handled by Proposition~\ref{tor:mask-degrees}.

For $d\ge4$, Theorem~\ref{tor:boundary} injects the dimension-$d$ group into the dimension-$(d-1)$ coatom groups. Induction proves both assertions. Every low-dimensional model and source was constructed over $\C$ and transported to the same $F$, so the induction never compares different cell function fields.
\end{proof}

\subsection{Type \texorpdfstring{$C$}{C}: the equivariant induction statement}
The remaining classical case needs more than local vanishing. Use signed ambient coordinates $\eps_i$ for the type-$C$ roots. A half-character $\eps_i$ need only belong to a finite covering torus; the actual coordinate characters below are roots $2\eps_i$ and $\eps_i+\eps_j$.

\begin{theorem}\label{tor:veronese-classification}
Let $N$ be an actual complex type-$C$ Schubert slice of dimension $d$. If $T_F(N)\ne0$ for some $F\supset\C$, then $d\ge2$ and $N$ is isomorphic, preserving its vertex and ambient torus action, to the scalar Veronese cone $V_d=\A^d/\{\pm1\}$, where the involution acts by simultaneous negation. Its coordinate axes have characters given by $d$ distinct ambient long roots. For this model, $T_F(N)=\Z/2$ for every $F\supset\C$.
\end{theorem}

The quantifier over every $F$ and the equivariance are part of the induction statement. An abstract analytic model, or a formal model without its characters, would not suffice in the next step.

\begin{lemma}\label{tor:veronese-base}
Theorem~\ref{tor:veronese-classification} holds for $d\le3$.
\end{lemma}
\begin{proof}
For $d\le1$, the puncture is empty or the spectrum of a field. For $d=2,3$, use the actual Levi reduction and Lemma~\ref{tor:products}. The single type-$A$ factors are covered by Proposition~\ref{tor:simply-laced-b}. The finite inputs leave only the $C_2$ pair $(s_1,s_1s_2s_1)$ and the $C_3$ pair with words $(Q_u,Q_w)=(121,121321)$.

Here is an actual equivariant model for both. In a symplectic space $E\oplus E^*$ with $\dim E=m$, the graph chart of $\LG(m,2m)$ consists of symmetric maps $S:E\longrightarrow E^*$. The Schubert incidence $Z_m=\{L:\dim(L\cap E)\ge m-1\}$ is the reduced rank-at-most-one symmetric-matrix locus in this chart. Choosing a hyperplane in $E$ and then its isotropic extension gives an irreducible $\PP^1$-bundle incidence, so $Z_m$ is integral of dimension $m$. Its top minimal word is $v=s_1\cdots s_m$. The full-flag inverse image has top $w=vw_J$ and lower endpoint $u=w_J$, where $J=\{1,\ldots,m-1\}$. On the opposite cell at $u$, the projection identifies the full-flag slice with the graph chart.

The finite map $x\longmapsto xx^{\mathsf t}$ has reduced image exactly the symmetric matrices of rank at most one. Its image ring is $\C[x_ix_j]$, the even invariant ring. For $m=2,3$, this gives the stated slices, since $123121=121321$. Their matrix characters are $\eps_i+\eps_j$, and their axis characters are $2\eps_i$. Conjugating a Levi changes these only by signed permutation and preserves root lengths. The next lemma supplies the torsion calculation.
\end{proof}

\subsection{A cover that cannot split on a normal subslice}
\begin{lemma}\label{tor:veronese-cover}
For $d\ge2$, the strict puncture of $V_d$ has \'etale fundamental group $\Z/2$ and Picard torsion $\Z/2$, over every $F\supset\C$. Its nontrivial double cover remains nonsplit on the puncture of every positive-dimensional normal local quotient of its strict vertex ring.
\end{lemma}
\begin{proof}
Let $A=A_F(V_d)$, and base change the finite scalar cover $F[x_1,\ldots,x_d]_{(x)}$ to $A$, obtaining $B$. Its closed fibre has one underlying point. Thus $B$ is local; it is strictly henselian and regular by finite henselian decomposition and ind-\'etale preservation of regularity. Before strict henselization, the invariant maximal ideal is generated by the quadratic monomials $x_ix_j$. After base change, $B/(x_1,\ldots,x_d)B$ is the separably closed residue field of $A$. Consequently the maximal ideals satisfy
\[
 \mathfrak n_B=(x_1,\ldots,x_d)B,
 \qquad \mathfrak m_AB=\mathfrak n_B^2,
\]
where the second identity uses the same quadratic generators.

The puncture of $B$ is a connected \'etale double cover of the puncture of $A$. Every finite \'etale cover of the regular strict puncture of $B$ extends to $\Spec B$ by purity \cite[Tag~0BMA, assertion~(1)]{The26}, and splits there. Hence the downstairs fundamental group is $\Z/2$. Normal Hartogs extension and strict henselian roots of units identify $H^1(U,\mu_m)$ with $\Pic(U)[m]$. Indeed, $\Gamma(U,\OO_U)^*=A^*$, and each unit has an $m$th root, so the unit term in the Kummer exact sequence vanishes. Thus $\Pic(U)[m]=\Hom(\Z/2,\mu_m)$: there is exactly one nonzero torsion class, of order two, corresponding to the scalar double cover.

Now let $A\twoheadrightarrow R$ be a positive-dimensional normal local domain. If the cover splits on its puncture, choose a generic branch to obtain $B\otimes_A R\longrightarrow\Frac R$. Integrality and normality put its image in $R$, giving an $A$-algebra map $B\longrightarrow R$. This map is local: the inverse image of $\mathfrak m_R$ lies over $\mathfrak m_A$, and $B$ has a unique prime over $\mathfrak m_A$, namely $\mathfrak n_B$. The quadratic-ideal identity and surjectivity of $A\longrightarrow R$ now imply $\mathfrak m_R=\mathfrak m_AR\subset\mathfrak m_R^2$. Nakayama's lemma contradicts the positive dimension of $R$. This also covers dimension one: the argument propagates a nonsplit cover, not nonzero Picard torsion on a curve.
\end{proof}

\subsection{Propagation along coatoms}
\begin{lemma}\label{tor:coatom}
Assume Theorem~\ref{tor:veronese-classification} in every dimension less than $d\ge4$. If a dimension-$d$ slice has nonzero $T_F$, every coatom slice is an equivariant scalar Veronese cone, and every proper upward root curve at its vertex has long-root character.
\end{lemma}
\begin{proof}
Theorem~\ref{tor:boundary} and induction inject $T_F$ into a sum of groups each equal to zero or $\Z/2$. Thus a nonzero class has order two. Its nontrivial Kummer double cover is nontrivial on at least one coatom. By induction that coatom is an equivariant Veronese cone, and by Lemma~\ref{tor:veronese-cover} the restricted cover is its scalar cover.

Suppose another coatom shares a lower label $z>u$ with this one. The actual shared-field quotient $A_F(u,z)$ is positive-dimensional and normal, so the restricted cover is nonsplit there by Lemma~\ref{tor:veronese-cover}. If its Picard class on the second coatom were zero, normal Hartogs extension and strict roots of units would make the cover split on the whole second-coatom puncture, a contradiction.

The graph of coatoms joined when they share a lower label $z>u$ is connected. Indeed, the comparability graph of $(u,w)$ is connected; choose a coatom above each vertex. The coatoms chosen above comparable vertices share the smaller vertex. Nontriviality therefore propagates to every coatom, and induction identifies them all equivariantly as Veronese cones.

Every proper upward root curve lies in one of these coatom slices. Its effective torus is the full diagonal torus of the scalar Veronese cone, whose one-dimensional orbit closures through the vertex are its coordinate axes. These axes have long-root characters. This argument makes no assumption about a possible direct edge from $u$ to $w$.
\end{proof}

\begin{lemma}\label{tor:short-edge}
In the situation of Lemma~\ref{tor:coatom}, a direct root edge from $u$ to $w$ cannot be short.
\end{lemma}
\begin{proof}
A reflection changes length by an odd integer, so if its gap is $d\ge4$, then $d\ge5$. The short positive roots in $C_n$ are the $D_n$ roots $e_i\pm e_j$ for $i<j$. Let $\ell_D^*(x)$ count the roots in this set inverted by a signed permutation $x$. Then $\ell_C(x)=\ell_D^*(x)+\operatorname{neg}(x)$, where $\operatorname{neg}(x)$ counts negative entries. A short reflection changes this count by at most two, and therefore $\ell_D^*(w)-\ell_D^*(u)\ge3$.

The endpoints lie in the same signed-parity coset. For the even coset, use $W(D_n)$ itself; for the odd coset, multiply on the left by the sign change $\tau$ in the last coordinate. This $\tau$ interchanges $e_i-e_n$ and $e_i+e_n$ and fixes the other positive $D_n$ roots. It preserves that positive system and $\ell_D^*$. For a positive root $\alpha$, the reflection criterion for Bruhat order says $x<s_\alpha x$ if and only if $x^{-1}\alpha$ is positive \cite[Chapter~2]{BB05}. The positive $D_n$ roots are precisely the positive short $C_n$ roots, and signed permutations preserve the set of short roots. Thus reflection-ascent directions agree in the two systems. In the odd coset, conjugate the reflection by $\tau$; the same criterion proves agreement after transport.

The transported endpoints are consequently comparable in $D_n$, with length gap at least three. Choose the first cover of a saturated $D_n$ chain between them. Its upper endpoint is strictly between the two endpoints. Every step transports back to a short-root reflection ascent in $C_n$, so this first endpoint gives a proper upward short-root curve between $u$ and $w$, contrary to Lemma~\ref{tor:coatom}. In ranks two and three, interpret $D_2=A_1\times A_1$ and $D_3=A_3$; rank one has no short roots. Only this order-preserving direction is used, not equality of induced Bruhat orders.
\end{proof}

\subsection{The actual toric ring}
\begin{proof}[Proof of Theorem~\ref{tor:veronese-classification}]
Lemma~\ref{tor:veronese-base} gives the low-dimensional bases. Proceed by dimension, simultaneously over all ranks and fields. For $d\ge4$, a nonzero class makes every upward curve long by Lemmas~\ref{tor:coatom}--\ref{tor:short-edge}. If there are $b$ such curves, Deodhar's upward inequality gives $b\ge d$ \cite[Proposition~2.3]{CK03}: there are at least $\ell(w)-\ell(u)$ reflections $s_\alpha$ with $u<s_\alpha u\le w$. The corresponding projective root curves lie in $X_w$, and their affine pieces through $u$ in the opposite cell are actual root lines in $N$.

The coordinate characters of this attractive chart lie in a pointed root cone, so both signs of a long axis cannot occur. Distinct long-root curves therefore have independent signed axis directions. Their characters belong to the effective weight span of $\C[N]$, giving $d\le b\le\rank T_{\mathrm{eff}}\le d$.
For the last inequality, choose a point where every nonzero root-coordinate generator is nonzero; its stabilizer is the action kernel. Equality proves that $T_{\mathrm{eff}}$ has a dense open orbit. Since $N$ is normal, it is an actual normal affine toric variety, not merely a toric tangent cone. We use the standard semigroup and orbit--face descriptions \cite[Chapters~1 and 3]{CLS11}.

Discard opposite-cell root coordinates that vanish identically on $N$, and normalize each remaining coordinate to have value one at a complex point of the dense orbit. On that orbit, each coordinate is its weight character. Restriction to the orbit is injective, and these coordinates generate $\C[N]$. Thus $\C[N]=\C[S]$, where $S$ is the semigroup generated by the actual root characters, and $M=\operatorname{gp}S$ is its actual character lattice.

Let $\sigma^\vee=\mathbf R_{\ge0}S$ be the character cone. For a face $\theta$ of this cone, quotienting by all monomials whose weights lie outside $\theta$ gives $\C[S\cap\theta]$. Consequently its one-dimensional faces correspond precisely to invariant irreducible curves through the vertex. Viewed in $G/B$, these are the upward root curves already counted. Their rays are exactly the $d$ independent long axes given by the preceding dimension equality, so the character cone is their positive orthant. The only ambient root-coordinate characters in that orthant are $2\eps_i$ and $\eps_i+\eps_j$.

Each $2\eps_i$ occurs, by restricting its root coordinate to the actual root line. Join $i$ to $j$ if a cross weight $\eps_i+\eps_j$ occurs, and let $J_a$ be the connected components of this graph. The lattice contains $2\Z^d$. Modulo that sublattice, edge vectors span the even-total subspace on each connected component, as follows by using paths in a spanning tree. Therefore
\[
 M=\left\{b\in\Z^d:
   \sum_{i\in J_a}b_i\equiv0\pmod2
   \text{ for every }a\right\}.
\]
Normality of $\C[S]$ implies $S=M\cap\Z_{\ge0}^d$. More explicitly, every $b$ in this intersection has a positive multiple in $S$: because all $2\eps_i$ occur, already $2b\in S$. The character of $b$ is integral over $\C[S]$ in its fraction field, hence belongs to $\C[S]$. This uses normality in $M$, without replacing $M$ by a larger ambient saturation.

This description of $M$ identifies the actual algebra, with its characters, as the tensor product of scalar Veronese invariant rings on the $J_a$. If two or more factors occur, Lemma~\ref{tor:products} gives $T_F(N)=0$, a contradiction. There is one factor, so $N=V_d$ equivariantly. Lemma~\ref{tor:veronese-cover} computes its group over every $F$ and completes the induction.
\end{proof}

\paragraph{A four-dimensional model.}
The rank-at-most-one symmetric $4\times4$ matrices form $\Spec\C[x_ix_j:1\le i\le j\le4]=V_4$. Its character lattice is $\{b\in\Z^4:\sum_i b_i\text{ is even}\}$, whereas its character cone is the full positive orthant. The four axis coordinates have weights $2\eps_i$, and the six off-diagonal coordinates have weights $\eps_i+\eps_j$. Setting all coordinates involving a fixed index equal to zero gives a closed $V_3$. Its scalar double cover remains nonsplit on each such punctured face, as required by the propagation argument. The distinction between the even-sum lattice and $\Z^4$ is essential: saturation in $\Z^4$ would give $\A^4$ and erase the order-two obstruction.

\begin{corollary}\label{tor:c-nonmaximal}
Let $X$ be a type-$C$ full-flag Schubert variety and let $Y\subseteq\Sing(X)$ be a Schubert subvariety that is not an irreducible component of $\Sing(X)$. The strict punctured Picard group of $\OO_{X,\eta_Y}$ has no torsion.
\end{corollary}
\begin{proof}
The only possible nonzero case in Theorem~\ref{tor:veronese-classification} is a slice isolated at its vertex. By the actual product chart of Subsection~\ref{tor:charts}, its label would then be the generic point of a maximal singular component. It cannot be a nonmaximal singular label.
\end{proof}

\Needspace{5\baselineskip}
\subsection{Exceptional types and the conclusion}
\begin{proposition}\label{tor:exceptional}
In $F_4$, $T_F(N)=0$ whenever $\dim N\ge4$. In $G_2$, it vanishes whenever $\dim N\ge3$. Every remaining nonzero group needed here occurs only at the generic point of a maximal singular component.
\end{proposition}
\begin{proof}
For $F_4$ in dimensions at most three, the actual Levi factors are of type $A$, $B$, $C$, or products. The preceding results leave only isolated scalar Veronese models as possible nonzero groups. Here a two-dimensional type-$B$ slice reduces again to rank at most two; the rank-two root system $B_2$ is isomorphic to $C_2$ after interchanging the node labels, and is covered by the surface calculation and graph model above. This also treats rank-two Levi factors inside $B_3$. Dimension four is Proposition~\ref{tor:certificates}, including the empty-column cases. Higher dimensions follow inductively from Theorem~\ref{tor:boundary}.

In $G_2$, reduced words alternate and have length at most six. Deleting an end letter gives length one less. Deleting an interior letter creates exactly one repeated adjacent letter; Demazure idempotence removes the repetition and leaves an alternating word of length two less. Thus single-deletion centres have gap one or two, and those of codimension at least two have gap exactly two. For gap at least three the column set is empty, so Proposition~\ref{tor:mask-degrees} gives vanishing. Normality implies that a nonmaximal singular label has local dimension at least three: in dimension two, normality leaves only an isolated singular vertex in the attractive slice. This covers every required $G_2$ label.
\end{proof}

\begin{theorem}\label{tor:torsion}
For every finite crystallographic root system and every Schubert variety $X\subset G/P$, the ordinary localization map
\[
 \Cl(\OO_{X,x})\longrightarrow
 \bigoplus_{\substack{Z\in\Irr(\Sing X)\\x\in Z}}
 \Cl(\OO_{X,\eta_Z})
\]
is injective on the torsion subgroup, for every scheme point $x$.
\end{theorem}
\begin{proof}
For an irreducible root system, the preceding propositions and Corollary~\ref{tor:c-nonmaximal} kill strict punctured Picard torsion at every nonmaximal singular label. For a reducible root system, the actual slices are products. At least two active positive-dimensional factors are treated by Lemma~\ref{tor:products}; a single active factor is already covered. Take $F=\C(C_u)$ in the generic-point description of Subsection~\ref{tor:charts}, and apply Proposition~\ref{tor:ordinary} to obtain detection at every ordinary point of a full-flag Schubert variety. Proposition~\ref{tor:partial} descends this assertion to all partial flags. This passage uses neither invariance of the entire class group nor any classification in a parabolic Bruhat quotient interval.
\end{proof}

Theorem~\ref{tor:torsion} proves Theorem~1.3.

